\subsection{ICache 集成}

需适配 pre-IF 级与 IF 级之间的数据流和控制流，并修改 AXI 总线接口模块以支持 Burst 传输。

\subsubsection{ICache 实例化与接口连接}

\begin{lstlisting}[language=Verilog, caption=ICache 实例化]
// ICache 模块实例化
cache #(
    .NUM_WAYS(2)
) u_icache (
    .clk(clk),
    .resetn(resetn),
    
    // CPU 侧接口
    .valid(inst_sram_req),
    .op(1'b0),                    // ICache 只读不写
    .index(inst_sram_addr[11:4]), // 虚地址 [11:4] 作为 Index
    .tag(inst_paddr[31:12]),      // 物理地址 [31:12] 作为 Tag
    .offset(inst_sram_addr[3:0]),
    .wstrb(4'b0),
    .wdata(32'b0),
    .cacheable(inst_mat != 2'b10),  // MAT != 2'b10 表示可缓存
    // ... 略
    .cacop(if_inst_icache_op),      // CACOP 操作 ICache
    .cacop_mode(if_cacop_mode),
    
    // AXI 侧接口
    .rd_req(icache_rd_req),
    .rd_type(icache_rd_type),
    // ... 略
    .rd_addr(icache_rd_addr),
);

// ICache 不会发起写请求，wr_rdy 恒为 1
assign icache_wr_rdy = 1'b1;
\end{lstlisting}

说明：
\begin{itemize}
    \item \textbf{Index 使用虚地址}：VIPT 设计，使用虚地址 [11:4] 索引 Cache
    \item \textbf{Tag 使用物理地址}：Tag 比较使用 TLB 转换后的物理地址 [31:12]
    \item \textbf{可缓存判断}：根据 MAT（Memory Access Type）域判断是否可缓存
    \item \textbf{CACOP}：CACOP 指令可以操作 ICache 进行初始化或失效
\end{itemize}

\subsubsection{AXI Burst 传输支持}

AXI 总线接口模块需要支持 Burst 传输模式：

\begin{lstlisting}[language=Verilog, caption=AXI Bridge Burst 传输]
// 读通道状态机
localparam R_IDLE = 3'b001;
localparam R_INST = 3'b010;
localparam R_DATA = 3'b100;

reg [2:0] r_state, r_next_state;
reg [3:0] inst_burst_cnt, data_burst_cnt;  // Burst 计数器

// 读地址通道
always @(*) begin
    if (r_state == R_IDLE) begin
        if (icache_rd_req) begin
            arid = 4'b0001;  // ICache ID
            araddr = icache_rd_addr;
            arlen = (icache_rd_type == 3'b100) ? 8'd3 : 8'd0;  // 4 拍或 1 拍
            arsize = 3'b010;  // 4 字节
            arburst = 2'b01;  // INCR
            arvalid = 1'b1;
        end else if (dcache_rd_req) begin
            arid = 4'b0010;  // DCache ID
            // ... 类似 ICache ，略
        end
    end
end

// 读数据通道
always @(posedge clk) begin
    if (!resetn) begin
        inst_burst_cnt <= 4'b0;
        data_burst_cnt <= 4'b0;
    end else begin
        if (r_state == R_INST && rvalid && rready) begin
            inst_burst_cnt <= inst_burst_cnt + 1;
        end
        if (r_state == R_DATA && rvalid && rready) begin
            data_burst_cnt <= data_burst_cnt + 1;
        end
    end
end

// 返回信号
assign icache_ret_valid = (r_state == R_INST) && rvalid;
assign icache_ret_last = (inst_burst_cnt == 4'd3) && rlast;
assign icache_ret_data = rdata;
// DCache 返回信号类似，略
\end{lstlisting}

关键设计说明：
\begin{itemize}
    \item \textbf{arlen 设置}：Cache 行访问 arlen=3（4 拍），字访问 arlen=0（1 拍）
    \item \textbf{Burst 计数}：通过记录已接收的数据拍数，判断 Burst 是否完成
\end{itemize}

\subsubsection{CACOP 指令对 ICache 的操作}

\begin{lstlisting}[language=Verilog, caption=CACOP 操作 ICache 的流水线控制]
// ID 级识别 CACOP 指令
wire id_inst_cacop = /* CACOP 指令译码 */;
wire id_inst_icache_op = id_inst_cacop && (cacop_code == 5'h00 || 
                                           cacop_code == 5'h01 ||
                                           cacop_code == 5'h02);

// IF 级传递 CACOP 信息
always @(posedge clk) begin
    if (rst || flush) begin
        if_inst_icache_op <= 1'b0;
        if_cacop_mode <= 2'b00;
    end else if (if_validin && if_allowin) begin
        if_inst_icache_op <= id_inst_icache_op;
        if_cacop_mode <= id_cacop_mode;
    end
end

// CACOP 操作 ICache 后需清空流水线
wire icache_cacop_flush = wb_inst_icache_op && wb_validout;

assign flush = wb_ex || ertn_flush || br_taken || icache_cacop_flush;
\end{lstlisting}

CACOP 操作 ICache 时，由于会修改 Cache 内容，需要清空流水线并重新取指，以保证取指一致性。\verb|cacop_mode| 决定具体操作：
\begin{itemize}
    \item \textbf{mode 0}：Index Invalid（索引失效）
    \item \textbf{mode 1}：Hit Invalid（命中失效）
    \item \textbf{mode 2}：Clear（命中失效）
\end{itemize}
